\chapter{Probabilités conditionnelles et arbres}
\section*{Exercice 22.1 -- Du tableau à la probabilité}
On choisit uniformément une personne dans le tableau du casque du chapitre 20 : A avec casque $90$, A sans casque $30$, B avec casque $40$, B sans casque $40$. Calculer $P(A)$, $P_A(C)$, $P_C(A)$ et $P(A\cap C)$.

\section*{Exercice 22.2 -- Deux chemins}
Un objet provient de l'atelier A avec probabilité $0{,}6$, sinon de B. La proportion de défauts est $2\,\%$ dans A et $5\,\%$ dans B. Construire un arbre. Calculer la probabilité d'un défaut et la part des objets de A parmi les objets défectueux.

\section*{Exercice 22.3 -- Un diagnostic fictif}
Dans un modèle fictif, $1\,\%$ d'une population est malade. Le test est positif chez $90\,\%$ des malades et chez $5\,\%$ des non-malades. Construire un tableau de proportions ramenées à $10\,000$ personnes. Un positif a-t-il $90\,\%$ de probabilité d'être malade ?

\section*{Exercice 22.4 -- Sensibilité et spécificité}
Dans le modèle précédent, donner sensibilité, spécificité, taux de faux négatifs parmi les malades et taux de faux positifs parmi les non-malades. Donner aussi la probabilité d'un faux positif dans toute la population.

\section*{Exercice 22.5 -- Indépendance (facultatif)}
Pour deux événements de probabilités non nulles, on a $P(A)=0{,}4$, $P(B)=0{,}5$ et $P(A\cap B)=0{,}2$. Sont-ils indépendants ? Sont-ils incompatibles ? Justifier par les définitions.

\section*{Exercice 22.6 -- Pourquoi la multiplication ?}
On suppose $P(A)>0$. À partir de la définition d'une probabilité conditionnelle, démontrer la règle du produit le long d'un chemin. Expliquer ce qui change si $P(A)=0$.
